\section{Результаты моделирования}

\subsection{Параметры модели}

Численное моделирование выполнено для радиального подшипника скольжения геометрии~\varGeomLetter{} ($R = \varRmm$~мм, $c = \varCmkm$~мкм, $L = \varLmm$~мм) с \varDimpleTypeGenitive{} микрорельефом. Диаметр вала составляет $D = 2R = 50$~мм. Данная геометрия является наименьшей из рассматриваемых (геометрия~A: $R = 35$~мм, геометрия~B: $R = 50$~мм), что позволяет оценить влияние масштабного фактора на эффективность текстурирования. Исходные данные для расчёта представлены в таблице~\ref{tab:input_data}.

\begin{table}[H]
\small
\caption{Исходные данные для расчёта}
\label{tab:input_data}
\begin{tabularx}{\textwidth}{|X|c|c|}
\hline
\centering\textbf{Параметр} & \textbf{Значение} & \textbf{Размерность} \\
\hline
Радиус вала $R$ & \varRmm & мм \\
\hline
Радиальный зазор $c$ & \varCmkm & мкм \\
\hline
Длина подшипника $L$ & \varLmm & мм \\
\hline
Относительная длина $L/D$ & \varLD & -- \\
\hline
Полуось по оси $Z$, $a$ & \varDimpleAmm & мм \\
\hline
Полуось по оси $\varphi$, $b$ & \varDimpleBmm & мм \\
\hline
Глубина углубления $h_p$ & \varDimpleHpmkm & мкм \\
\hline
Относительная глубина $h_p/c$ & \varDimpleHpRatio & -- \\
\hline
Число углублений по окружности $N_\varphi$ & \varNphi & шт. \\
\hline
Число рядов по длине $N_Z$ & \varNZ & шт. \\
\hline
Область текстурирования & \varPhiRange & -- \\
\hline
Частота вращения $n$ & \varN & об/мин \\
\hline
Вязкость смазки $\eta$ & \varEta & Па$\cdot$с \\
\hline
Диапазон эксцентриситета $\varepsilon$ & 0{,}05--0{,}80 & -- \\
\hline
Размер сетки & $500 \times 500$ & -- \\
\hline
Параметр релаксации $\omega$ & 1{,}5 & -- \\
\hline
Критерий сходимости $\delta$ & $< 10^{-5}$ & -- \\
\hline
Макс. число итераций & 50000 & -- \\
\hline
\end{tabularx}
\end{table}

Линейная скорость поверхности вала при заданной частоте вращения равна $U = \pi D n / 60 = \pi \cdot 0{,}050 \cdot 2980 / 60 \approx 7{,}80$~м/с. По сравнению с подшипниками геометрий~A и~B линейная скорость существенно ниже вследствие меньшего диаметра вала, что приводит к пропорциональному снижению абсолютных значений нагрузочной способности и расхода смазки. Расчёт проведён для 15 значений эксцентриситета в указанном диапазоне.

\subsection{Распределение зазора и давления}

На рисунке~\ref{fig:pressure_2d} представлено распределение безразмерного давления $P(\varphi)$ в среднем сечении ($Z = 0$) при $\varepsilon = 0{,}6$ для гладкого и текстурированного подшипников.

\begin{figure}[H]
\centering
\includegraphics[width=\textwidth]{figures/pressure_2d.png}
\caption{Распределение давления $P(\varphi)$ при $Z=0$ для гладкого и текстурированного подшипников ($\varepsilon = 0{,}6$)}
\label{fig:pressure_2d}
\end{figure}

Кривая давления гладкого подшипника имеет характерную колоколообразную форму с максимумом вблизи области минимального зазора ($\varphi \approx \pi$). В расширяющейся части зазора давление обращается в нуль вследствие кавитации, а область положительного давления охватывает приблизительно половину окружности. Введение цилиндрического микрорельефа приводит к появлению выраженных локальных возмущений на кривой давления. Характерной особенностью цилиндрического профиля является наличие резких перепадов давления на границах углублений, обусловленных скачкообразным изменением зазора. В отличие от эллипсоидального профиля, где глубина меняется плавно, цилиндрическое углубление создаёт ступенчатый переход, что формирует более интенсивные пики давления на входных кромках. Для подшипника наименьшей геометрии~C с радиальным зазором $c = 30$~мкм относительная глубина углублений $h_p/c = 0{,}3$ выше, чем при геометриях~A и~B ($h_p/c = 0{,}2$), что усиливает микрогидродинамический эффект. Общий уровень давления текстурированного подшипника существенно превышает гладкий аналог, обеспечивая значительный прирост несущей способности.

На рисунке~\ref{fig:clearance_2d} изображён профиль зазора $H(\varphi)$ в среднем сечении при $\varepsilon = 0{,}6$.

\begin{figure}[H]
\centering
\includegraphics[width=\textwidth]{figures/clearance_2d.png}
\caption{Распределение зазора $H(\varphi)$ при $Z=0$ для подшипника с \varDimpleTypeGenitive{} углублениями ($\varepsilon = 0{,}6$)}
\label{fig:clearance_2d}
\end{figure}

Профиль зазора следует закону $H = 1 + \varepsilon\cos\varphi$, на фоне которого видны ступенчатые увеличения, характерные для цилиндрических углублений. Глубина каждого углубления в безразмерных единицах составляет $\Delta H_{\max} = h_p/c = 0{,}3$. Переходы на кромках углублений имеют вертикальный характер, что отличает цилиндрический профиль от плавных эллипсоидальных форм.

Трёхмерные распределения зазора и давления при $\varepsilon = 0{,}6$ показаны на рисунках~\ref{fig:clearance_3d} и~\ref{fig:pressure_3d}.

\begin{figure}[H]
\centering
\includegraphics[width=\textwidth]{figures/clearance_3d.png}
\caption{Трёхмерное распределение зазора $H(\varphi, Z)$ ($\varepsilon = 0{,}6$)}
\label{fig:clearance_3d}
\end{figure}

Рисунок~\ref{fig:clearance_3d} демонстрирует регулярный массив цилиндрических углублений: 8 по окружности и 11 рядов по длине подшипника, размещённых в зоне $90^\circ - 270^\circ$.

\begin{figure}[H]
\centering
\includegraphics[width=\textwidth]{figures/pressure_3d.png}
\caption{Трёхмерное распределение давления $P(\varphi, Z)$ ($\varepsilon = 0{,}6$)}
\label{fig:pressure_3d}
\end{figure}

Из рисунка~\ref{fig:pressure_3d} видно, что каждое цилиндрическое углубление формирует резкий подъём давления на входной кромке, где поток сталкивается со ступенчатым расширением зазора. Этот микрогидродинамический эффект при цилиндрическом профиле проявляется более отчётливо, чем при плавных формах, поскольку резкие границы углубления создают выраженные локальные градиенты давления. На выходной кромке наблюдается спад давления, однако кавитационное ограничение ($P \geq 0$) обеспечивает положительный суммарный вклад каждого углубления в несущую способность.

При рассмотрении трёхмерных распределений следует выделить ряд характерных закономерностей. Максимальное давление формируется в среднем сечении ($Z = 0$), а к торцам подшипника ($Z = \pm 1$) давление спадает до нуля из-за торцевых утечек. Для подшипника геометрии~C с малой длиной ($L = 40$~мм) торцевые утечки играют более заметную роль, чем для крупных подшипников. Микрорельеф оказывает наибольшее влияние в центральной области, где фоновое давление достигает максимума. Вблизи торцов эффект текстурирования ослабевает вследствие низкого уровня окружающего давления.

\subsection{Сравнительный анализ интегральных показателей}

В таблицах~\ref{tab:smooth_results} и~\ref{tab:textured_results} приведены интегральные характеристики гладкого подшипника и подшипника с \varDimpleTypeGenitive{} углублениями для характерных значений эксцентриситета.

\begin{table}[H]
\small
\centering
\caption{Характеристики гладкого подшипника}
\label{tab:smooth_results}
\begin{tabular}{|c|c|c|c|}
\hline
$\boldsymbol{\varepsilon}$ & $\boldsymbol{F}$\textbf{, Н} & $\boldsymbol{\mu}$ & $\boldsymbol{Q}$\textbf{, л/с} \\
\hline
0{,}05 & 187{,}7 & 0{,}1925 & 0{,}0734 \\
\hline
0{,}16 & 604{,}4 & 0{,}0606 & 0{,}0733 \\
\hline
0{,}26 & 1069{,}3 & 0{,}0353 & 0{,}0731 \\
\hline
0{,}37 & 1635{,}4 & 0{,}0241 & 0{,}0729 \\
\hline
0{,}48 & 2377{,}2 & 0{,}0177 & 0{,}0726 \\
\hline
0{,}59 & 3445{,}8 & 0{,}0134 & 0{,}0723 \\
\hline
0{,}69 & 5223{,}6 & 0{,}0101 & 0{,}0720 \\
\hline
0{,}80 & 8836{,}8 & 0{,}0073 & 0{,}0718 \\
\hline
\end{tabular}
\end{table}

\begin{table}[H]
\small
\centering
\caption{Характеристики подшипника с \varDimpleTypeGenitive{} углублениями}
\label{tab:textured_results}
\begin{tabular}{|c|c|c|c|}
\hline
$\boldsymbol{\varepsilon}$ & $\boldsymbol{F}$\textbf{, Н} & $\boldsymbol{\mu}$ & $\boldsymbol{Q}$\textbf{, л/с} \\
\hline
0{,}05 & 1083{,}2 & 0{,}0321 & 0{,}0781 \\
\hline
0{,}16 & 1464{,}2 & 0{,}0240 & 0{,}0780 \\
\hline
0{,}26 & 2019{,}8 & 0{,}0178 & 0{,}0778 \\
\hline
0{,}37 & 2819{,}6 & 0{,}0133 & 0{,}0775 \\
\hline
0{,}48 & 4008{,}6 & 0{,}0099 & 0{,}0772 \\
\hline
0{,}59 & 5971{,}0 & 0{,}0072 & 0{,}0769 \\
\hline
0{,}69 & 9396{,}2 & 0{,}0051 & 0{,}0766 \\
\hline
0{,}80 & 16090{,}4 & 0{,}0035 & 0{,}0763 \\
\hline
\end{tabular}
\end{table}

Абсолютные значения нагрузочной способности подшипника геометрии~C существенно ниже, чем для крупных геометрий~A и~B, что обусловлено меньшими размерами вала и, соответственно, меньшей площадью несущей поверхности и более низкой линейной скоростью. Однако относительный эффект цилиндрического текстурирования проявляется значительно ярче: при $\varepsilon = 0{,}80$ нагрузочная способность текстурированного подшипника (16090{,}4~Н) превосходит гладкий аналог (8836{,}8~Н) на 82{,}1\%, что существенно выше, чем для геометрий~A и~B. Повышенная относительная глубина углублений ($h_p/c = 0{,}3$ против $0{,}2$) при меньшем зазоре усиливает микрогидродинамический эффект на резких кромках цилиндрических углублений.

Коэффициент трения $\mu$ уменьшается с ростом $\varepsilon$ для обоих вариантов, причём текстурированный подшипник имеет существенно более низкие значения $\mu$ во всём диапазоне. Расход смазки $Q$ практически не зависит от эксцентриситета и несколько выше у текстурированного подшипника вследствие увеличения эффективного объёма смазочного слоя в области углублений.

Количественное сопоставление характеристик при $\varepsilon = 0{,}6$ выполнено в таблице~\ref{tab:comparison}.

\begin{table}[H]
\small
\centering
\caption{Сравнение характеристик при $\varepsilon = 0{,}6$}
\label{tab:comparison}
\begin{tabular}{|l|c|c|c|}
\hline
\textbf{Характеристика} & \textbf{Гладкий} & \textbf{С углублениями} & $\boldsymbol{\Delta}$\textbf{, \%} \\
\hline
Нагрузка $F$, Н & 3636{,}8 & 6312{,}1 & $+73{,}56$ \\
\hline
Коэфф. трения $\mu$ & 0{,}0129 & 0{,}0069 & $-46{,}76$ \\
\hline
Расход $Q$, л/с & 0{,}0722 & 0{,}0769 & $+6{,}42$ \\
\hline
\end{tabular}
\end{table}

Применение цилиндрических углублений на подшипнике наименьшей геометрии~C обеспечивает прирост несущей способности на 73{,}56\%, что является наиболее значительным эффектом среди всех трёх рассмотренных геометрий с цилиндрическим профилем. Коэффициент трения снижается на 46{,}76\% --- столь существенное уменьшение обусловлено одновременным действием двух факторов: ступенчатого профиля углублений, создающего интенсивные микрогидродинамические эффекты, и повышенной относительной глубины $h_p/c = 0{,}3$. Увеличение расхода смазки на 6{,}42\% несколько выше, чем для крупных геометрий, что связано с более значительным вкладом углублений в общий объём смазочного зазора при малых абсолютных размерах подшипника.

На рисунках~\ref{fig:load_vs_eps}--\ref{fig:flow_vs_eps} показаны зависимости интегральных характеристик от эксцентриситета для обоих вариантов подшипника.

\begin{figure}[H]
\centering
\includegraphics[width=\textwidth]{figures/load_vs_epsilon.png}
\caption{Нагрузочная способность $F(\varepsilon)$ для гладкого и текстурированного подшипников}
\label{fig:load_vs_eps}
\end{figure}

Зависимость $F(\varepsilon)$ текстурированного подшипника расположена значительно выше кривой гладкого подшипника на всём интервале эксцентриситетов. Разрыв между кривыми нарастает с увеличением $\varepsilon$: при $\varepsilon = 0{,}05$ разница составляет около 896~Н, тогда как при $\varepsilon = 0{,}80$ она превышает 7250~Н. Нелинейный рост обусловлен кубической зависимостью давления от обратной толщины зазора ($P \sim h^{-3}$), вследствие чего при малых зазорах эффект каждого углубления многократно усиливается. Примечательно, что даже при малых эксцентриситетах текстурированный подшипник геометрии~C обеспечивает нагрузочную способность, многократно превосходящую гладкий аналог.

\begin{figure}[H]
\centering
\includegraphics[width=\textwidth]{figures/friction_vs_epsilon.png}
\caption{Коэффициент трения $\mu(\varepsilon)$ для гладкого и текстурированного подшипников}
\label{fig:friction_vs_eps}
\end{figure}

Для текстурированного подшипника коэффициент трения ниже гладкого при всех значениях $\varepsilon$. Кривые монотонно убывают и асимптотически стремятся к нулю при $\varepsilon \to 1$. Максимальный абсолютный выигрыш приходится на малые эксцентриситеты: при $\varepsilon = 0{,}05$ разность $\Delta\mu = 0{,}160$, тогда как при $\varepsilon = 0{,}80$ она составляет лишь $\Delta\mu = 0{,}004$. Относительное снижение трения остаётся значительным на протяжении всего диапазона эксцентриситетов, что свидетельствует об устойчивости эффекта цилиндрического текстурирования для малоразмерных подшипников.

\begin{figure}[H]
\centering
\includegraphics[width=\textwidth]{figures/flow_vs_epsilon.png}
\caption{Расход смазки $Q(\varepsilon)$ для гладкого и текстурированного подшипников}
\label{fig:flow_vs_eps}
\end{figure}

Расход смазки текстурированного подшипника превышает расход гладкого. Разница находится в пределах 6--7\%, что заметно больше, чем для подшипников геометрий~A и~B (1{,}5--2\%). Это объясняется тем, что при малых абсолютных размерах подшипника объём углублений составляет более значительную долю от общего объёма смазочного слоя.

Проведённый анализ свидетельствует о том, что цилиндрическое текстурирование на подшипнике наименьшей геометрии~C позволяет достичь наибольшего относительного эффекта по нагрузочной способности и коэффициенту трения. Одновременное существенное повышение нагрузочной способности и уменьшение коэффициента трения при умеренном приросте расхода смазки делает данный тип микрорельефа перспективным для малоразмерных подшипниковых узлов.

\subsection*{Выводы по главе 3}
\phantomsection
\addcontentsline{toc}{subsection}{Выводы по главе 3}

Проведён расчёт рабочих характеристик радиального подшипника скольжения геометрии~C с \varDimpleTypeGenitive{} микрорельефом для 15 значений эксцентриситета. Построены поля давления и зазора, определены зависимости нагрузочной способности, коэффициента трения и расхода смазки от эксцентриситета. Сравнение при $\varepsilon = 0{,}6$ выявило прирост нагрузочной способности на $73{,}56$\%, снижение коэффициента трения на $46{,}76$\% и увеличение расхода на $6{,}42$\%, что свидетельствует о высокой эффективности цилиндрического текстурирования, особенно выраженной для подшипников малых размеров.
